\section{Introduction}
\label{sec:introduction}
Sparse spectral methods seek low-dimensional directions that preserve important variation while depending on a small number of variables. The generalized power method of \citet{journee2010gpower} is influential because it turns high-dimensional sparse-loading problems into thresholded convex objectives over a sphere or Stiefel manifold. The associated linear maximization oracle is particularly simple: normalization on a sphere and the polar factor on a Stiefel manifold. This combination produces low per-iteration cost and exposes a useful separation between support discovery and support-conditioned refitting.

The multi-component construction, however, contains three mathematical ambiguities that propagate into its interpretation and implementation. First, eliminating each loading column from the printed linear-minus-\(\ell_1\) primal gives a sum of thresholded norms, while the displayed reduced block objective is a sum of their squares. Those objectives coincide in the single-component case only because a nonnegative scalar and its square have the same maximizers; they are not equivalent after summing several components. Second, the printed unit-sphere loading domain excludes zero loading columns even though the positive-part reductions and accompanying text allow inactive components. Third, the published post-processing routine normalizes a loading column before applying the support mask, which generally destroys unit norm after masking. These issues change the estimator, feasible set, fixed points, and validity of subsequent convergence arguments.

This paper treats the formulation audit as the beginning rather than the end of the research problem. We define two coherent block estimators. The first, \emph{A-SUM}, is the exact score-space reduction of a normalized loading-ball primal and uses a sum of thresholded norms. The second, \emph{A-SQ}, uses a sum of squared thresholded norms and is justified by a distinct quadratic-lift primal. We then study two algorithms for A-SQ. The historical full-polar step is retained for reproducibility and is analyzed only to the strength supported by convex minorization. A separate safeguarded Riemannian method is introduced for whole-sequence analysis; it combines an ascent direction, a retraction or polar candidate, an Armijo test, and an explicit fallback rule.

The central dynamical thesis is two-phase. Before identification, the method crosses threshold boundaries and changes a discrete support and sign pattern. Under a strict margin, convergence implies that these changes stop after finitely many iterations. After identification, the objective is smooth on a fixed active region and the iteration becomes a spectral map. In the rank-one hard-threshold case this map is ordinary power iteration on the active covariance. In the block case with a common active operator it is orthogonal iteration. With component-specific supports, the correct object is instead the Fr\'echet derivative of a coupled polar map on the Stiefel tangent space. This distinction prevents an overly broad claim that every block instance reduces to one standard covariance eigenproblem.

The modern-AI application is chosen to test, rather than replace, this mathematics. Transformer activation matrices are wide, structured, and routinely analyzed through sparse autoencoders or native-neuron selection. We ask whether measurable threshold margins predict support stabilization, whether active eigengaps or Jacobian radii predict local slopes, and whether the resulting sparse components have reproducible causal effects. Reconstruction is not treated as evidence of interpretability on its own; the \Cref{sec:results} panels report top-token coherence against GPT-2 input embeddings (an SAEBench-style quality diagnostic \citep{karvonen2025saebench}), multi-probe ablation KL against a matched-norm random-direction null (an MIB-style causal diagnostic \citep{mueller2025mib,oneill2025compute}), and a head-to-head comparison with a gated-SAE-style tied-decoder dictionary of matching width \citep{rajamanoharan2024gated}.

\paragraph{Contributions.}
\begin{enumerate}[leftmargin=*]
  \item We derive the exact loading-ball reduction and the exact quadratic-lift representation, and exhibit a two-component counterexample proving that their block aggregations define different estimators.
  \item We give complete numerical conventions for inactive components, threshold ties, rank-deficient gradient matrices, nonunique polar factors, component permutations, and support-conditioned refitting.
  \item We separate the historical full-polar algorithm from a safeguarded method, prove full-polar monotonicity, and \emph{prove} last-iterate convergence of the safeguarded method via a sufficient-ascent lemma, a relative-error lemma, and the Kurdyka--\L{}ojasiewicz property, which holds because the soft-threshold objective and the algebraic Stiefel constraint are semi-algebraic.
  \item We prove conditional finite identification under strict margins, derive the rank-one fixed-support power rate (validated on a controlled fixture to within $5\%$ of the classical $\lambda_2/\lambda_1$ prediction, and to within $1\%$ of a corrected prediction that includes the soft-threshold shift $\gamma\,\sigma_1\sqrt{k}$), formulate the correct block tangent-Jacobian condition, and obtain a locally safe screening rule from certified distance radii; the screening rule is realized at the solver level by a dedicated screened operator and attains the full $(n-k)/n$ oracle-work ceiling with machine-precision objective agreement against the unscreened solve.
  \item We validate the theory on planted spikes and extend the pipeline to a GPT-2 small residual stream, a vision-transformer CLS stream, and the Jeffers Pitprops classical sparse-PCA benchmark \citep{jeffers1967pitprops}; FACET components match dense PCA on reconstruction $R^2$ while exceeding it on Hoyer sparsity, top-token coherence, and class-conditional specialisation, and avoid the ${\sim}8$-point $R^2$ cost that a tied-decoder sparse autoencoder incurs for its $\ell_1$-driven sparsity at matched dictionary width.
\end{enumerate}

\paragraph{Organization.}
\Cref{sec:related} positions the work; \Cref{sec:setup} fixes the models; \Cref{sec:method} gives the algorithms and refit; \Cref{sec:theory} develops the result hierarchy and \Cref{sec:extended-theory} contains the formal proofs of the whole-sequence convergence program; \Cref{sec:experiments,sec:results} present the controlled and modern-AI empirical pipeline.

\begin{figure}[t]
\centering
\begin{tikzpicture}[node distance=5mm and 10mm, every node/.style={font=\small}]
% left column, printed / seed-paper form (muted / neutral style)
\node[neutralbox, text width=0.42\linewidth] (p0) {\textbf{Printed (seed) form}\\[-0.5mm]{\scriptsize block $\ell_1$ with unit-norm loading columns}};
\node[neutralbox, text width=0.42\linewidth, below=of p0] (p1) {\textbf{Single reduced objective}\\[-0.5mm]{\scriptsize $\sum_j\norm{\Soft_{\gamma_j}(c_j)}_2^2$, sum and squared forms conflated}};
\node[neutralbox, text width=0.42\linewidth, below=of p1] (p2) {\textbf{Normalize-then-mask refit}\\[-0.5mm]{\scriptsize masked vector generally not unit norm}};
\node[neutralbox, text width=0.42\linewidth, below=of p2] (p3) {\textbf{First-order theory}\\[-0.5mm]{\scriptsize stationarity residual $\to 0$ only, no last-iterate / identification / local-rate claim}};

% right column, FACET corrected form (highlighted)
\node[successbox, text width=0.42\linewidth, right=of p0] (c0) {\textbf{Corrected form}\\[-0.5mm]{\scriptsize loading-ball primal with explicit zero-column (inactive) policy}};
\node[successbox, text width=0.42\linewidth, below=of c0] (c1) {\textbf{Distinct estimators}\\[-0.5mm]{\scriptsize $F_{\mathrm{sum}}(X)=\sum_j\norm{\Soft_{\gamma_j}(c_j)}_2$ and $F_{\mathrm{sq}}(X)=\sum_j\norm{\Soft_{\gamma_j}(c_j)}_2^2$, not interchangeable at $m\!\ge\!2$}};
\node[successbox, text width=0.42\linewidth, below=of c1] (c2) {\textbf{Mask-then-normalize refit}\\[-0.5mm]{\scriptsize $z\!=\!Mc/\norm{Mc}_2$, restores ball feasibility}};
\node[successbox, text width=0.42\linewidth, below=of c2] (c3) {\textbf{Safeguarded ascent + K\L{} theory}\\[-0.5mm]{\scriptsize whole-sequence convergence, finite support/sign identification, post-identification local rate}};

\draw[arrPrimary] (p0) -- node[above,font=\scriptsize]{correct} (c0);
\draw[arrPrimary] (p1) -- node[above,font=\scriptsize]{split}   (c1);
\draw[arrPrimary] (p2) -- node[above,font=\scriptsize]{reorder} (c2);
\draw[arrPrimary] (p3) -- node[above,font=\scriptsize]{strengthen} (c3);
\end{tikzpicture}
\caption{Printed-vs-corrected model map. The left column lists four defects of the seed-paper block formulation; the right column states the FACET correction that each defect motivates. The downstream theorems A1--A10 depend on all four corrections being in force simultaneously.}
\label{fig:a-printed-vs-corrected}
\end{figure}

\begin{figure}[t]
\centering
\begin{tikzpicture}[node distance=10mm and 12mm]
\node[primarybox,text width=3.6cm] (primal) {Normalized loading-ball primal\\$\max_{\|z_j\|\le1} c_j^\top z_j-\gamma_j\|z_j\|_1$};
\node[secondarybox,text width=3.5cm,right=of primal] (sum) {A-SUM\\$\sum_j\|\Soft_{\gamma_j}(c_j)\|_2$};
\node[primarybox,text width=3.6cm,below=of primal] (qlift) {Quadratic-lift primal\\$\max_z 2c_j^\top z-\|z\|^2-2\gamma_j\|z\|_1$};
\node[secondarybox,text width=3.5cm,right=of qlift] (sq) {A-SQ\\$\sum_j\|\Soft_{\gamma_j}(c_j)\|_2^2$};
\draw[arrPrimary] (primal)--node[above,font=\scriptsize]{exact elimination}(sum);
\draw[arrPrimary] (qlift)--node[above,font=\scriptsize]{exact elimination}(sq);
\draw[dashed,thick,draw=Warning] (sum)--node[right,font=\scriptsize,align=center]{not equivalent\\for $m>1$}(sq);
\end{tikzpicture}
\caption{The corrected model family. A-SUM and A-SQ have different exact primal interpretations and are not silently identified.}
\label{fig:a-objectives}
\end{figure}

\begin{table}[t]
\centering
\caption{Paper contract.}
\label{tab:paper-contract}
\small
\begin{tabularx}{0.97\linewidth}{@{}>{\bfseries}p{0.23\linewidth}X@{}}
\toprule
Central question & Which corrected block estimator is being optimized, and what deterministic dynamics follow before and after support identification? \\
Primary object & Separate sum-of-thresholded-norms and squared-threshold objectives on the Stiefel manifold. \\
Main theorem package & Exact reduction, safeguarded whole-sequence convergence, finite support/sign identification, and scoped local rates. \\
Modern-AI test & Sparse, nonredundant residual-stream or MLP-channel subspaces in open language models. \\
Kill condition & The corrected structure yields no distinct behavior, no support-stable phase, or no theorem beyond generic stationarity. \\
\bottomrule
\end{tabularx}
\end{table}

